% chapters/ch011_utilitarianism.tex

\section{Utilitarianism}

This chapter gives a mathematical argument for why utilitarianism can be a reasonable social welfare criterion. The key idea is that if utility is quasi-linear in money and society can implement balanced money transfers, then the policy problem separates into:
\begin{enumerate}
    \item \textbf{efficiency:} choose the allocation that maximizes the total utility generated by goods;
    \item \textbf{distribution:} use money transfers to redistribute that utility pie.
\end{enumerate}
Under these conditions, Pareto-efficient policies are exactly those built from utilitarian allocations.

\subsection{Mathematical Logic for Economic Statements}

Lecture 11 begins with a short review of mathematical logic because many economic results in this course take the form of conditional statements.

\begin{definition}{Conditional statement, contrapositive, and converse}{ch11-logic}
A \emph{conditional statement} is a proposition of the form
\[
A \to B,
\]
read as ``If $A$, then $B$.''

Its \emph{contrapositive} is
\[
\neg B \to \neg A,
\]
and its \emph{converse} is
\[
B \to A.
\]
\end{definition}

\begin{theorem}{Conditional statement and contrapositive are equivalent}{ch11-contrapositive}
The statements
\[
A \to B
\qquad \text{and} \qquad
\neg B \to \neg A
\]
are logically equivalent.
\end{theorem}

\begin{proof}
A conditional statement fails only in the case where $A$ is true and $B$ is false. But this is exactly the case where $\neg B$ is true and $\neg A$ is false, so the contrapositive also fails in precisely the same states of the world. Therefore the two statements always have the same truth value.
\end{proof}

\begin{commonmistake}
Do not confuse the \emph{contrapositive} with the \emph{converse}. In general,
\[
A \to B
\]
does \emph{not} imply
\[
B \to A.
\]
\end{commonmistake}

\begin{examplebox}[title={Worked Example: Using a contrapositive}]
Consider the statement:
\[
\text{If } x+y=1000,\ \text{then either } x\ge 500 \text{ or } y\ge 500.
\]
A direct proof must consider all cases with $x+y=1000$. It is easier to prove the contrapositive:
\[
\text{If } x<500 \text{ and } y<500,\ \text{then } x+y\neq 1000.
\]
But if both $x$ and $y$ are less than 500, then
\[
x+y<1000.
\]
So the contrapositive is true, and hence the original statement is also true.
\end{examplebox}

\begin{examtip}
In this lecture, several results are proved most cleanly by contraposition:
\[
A \to B
\quad \Longleftrightarrow \quad
\neg B \to \neg A.
\]
This is often easier than proving the original implication directly.
\end{examtip}

\subsection{Rawlsian and Utilitarian Social Welfare}

Lecture 10 introduced social welfare functions. Lecture 11 contrasts two central examples.

\begin{definition}{Rawlsian and utilitarian social welfare functions}{ch11-rawls-util}
For utility vector
\[
u=(u_1,\dots,u_n),
\]
two common social welfare functions are:
\begin{itemize}
    \item \textbf{Rawlsian:}
    \[
    W_R(u)=\min\{u_1,\dots,u_n\};
    \]
    \item \textbf{Utilitarian:}
    \[
    W_U(u)=\sum_{i=1}^n u_i.
    \]
\end{itemize}
\end{definition}

The Rawlsian criterion puts maximal emphasis on the worst-off individual, while the utilitarian criterion puts equal weight on increasing the sum of utilities.

\begin{policynote}
The lecture's goal is not to prove that utilitarianism is always morally correct. Rather, it gives a precise economic argument showing that utilitarian allocations have a special efficiency property once money transfers and quasi-linear utility are introduced.
\end{policynote}

\subsection{Social Choice with Money Transfers}

The lecture now enriches the social choice problem by allowing monetary transfers alongside allocations of goods.

\begin{definition}{Allocation, transfer scheme, and policy}{ch11-policy}
A \emph{transfer scheme} is a vector
\[
t=(t_1,\dots,t_n),
\]
where:
\begin{itemize}
    \item $t_i>0$ means individual $i$ receives money;
    \item $t_i<0$ means individual $i$ pays money.
\end{itemize}
A \emph{policy} is a pair
\[
(A,t),
\]
consisting of:
\begin{itemize}
    \item an allocation $A$ of goods;
    \item a monetary transfer scheme $t$.
\end{itemize}
\end{definition}

\begin{definition}{Balanced transfers}{ch11-balanced}
A transfer scheme is \emph{budget balanced} if
\[
\sum_{i=1}^n t_i = 0.
\]
Thus money is only redistributed across individuals; no outside surplus or deficit is created.
\end{definition}

\begin{theorem}{Why balanced transfers are needed}{ch11-balanced-needed}
If unrestricted transfer schemes are allowed, Pareto efficiency of policies becomes trivial: any policy can be Pareto improved upon by giving every individual additional money.
\end{theorem}

\begin{proof}
Suppose $(A,t)$ is any policy. If transfers need not balance, choose $\varepsilon>0$ and consider
\[
t'_i=t_i+\varepsilon
\qquad \text{for all } i.
\]
If utility is increasing in money, then every individual is strictly better off under $(A,t')$ than under $(A,t)$. Hence $(A,t)$ cannot be Pareto efficient. Therefore meaningful analysis requires balanced transfers.
\end{proof}

\begin{commonmistake}
Do not confuse an \emph{allocation} with a \emph{policy}. In this chapter:
\begin{itemize}
    \item an allocation concerns only goods;
    \item a policy concerns goods \emph{and} transfers.
\end{itemize}
Pareto efficiency is applied to policies, not just to allocations.
\end{commonmistake}

\subsection{Quasi-Linear Utility and Willingness to Pay}

The central assumption is that individuals value money directly and linearly.

\begin{definition}{Quasi-linear utility with transfers}{ch11-quasilinear}
Suppose that if allocation $A$ is chosen and individual $i$ receives transfer $t_i$, utility is
\[
U_i(A,t_i)=v_i(A)+t_i,
\]
where:
\begin{itemize}
    \item $v_i(A)$ is the utility derived from the allocation itself;
    \item $t_i$ is money, valued one-for-one in utility.
\end{itemize}
\end{definition}

This means money can be used to transfer utility across individuals in a direct way.

\begin{theorem}{Willingness to pay under quasi-linear utility}{ch11-wtp}
Suppose the status quo is allocation $A$, and the new allocation is $B$. If moving from $A$ to $B$ requires individual $i$ to pay $c_i$, then $i$ prefers the change if and only if
\[
c_i \le v_i(B)-v_i(A).
\]
Hence
\[
v_i(B)-v_i(A)
\]
is individual $i$'s maximum willingness to pay to move from $A$ to $B$.
\end{theorem}

\begin{proof}
Individual $i$ weakly prefers the new policy if
\[
U_i(B,-c_i)\ge U_i(A,0).
\]
Using quasi-linearity,
\[
v_i(B)-c_i \ge v_i(A),
\]
which is equivalent to
\[
c_i \le v_i(B)-v_i(A).
\]
\end{proof}

\begin{examplebox}[title={Worked Example: Movie ticket allocation with transfers}]
Suppose Joseph and Ivy both value a movie ticket, with willingness to pay
\[
WTP_J < WTP_I.
\]
If Joseph initially gets the ticket and no money changes hands, that allocation is Pareto efficient \emph{without} transfers: giving the ticket to Ivy would hurt Joseph.

Now allow money transfers. Give the ticket to Ivy and let Ivy pay Joseph an amount $x$ such that
\[
WTP_J \le x \le WTP_I.
\]
Then:
\begin{itemize}
    \item Joseph weakly prefers losing the ticket and receiving $x$;
    \item Ivy weakly prefers getting the ticket and paying $x$;
    \item at least one of them is strictly better off unless the inequalities are equalities at both ends.
\end{itemize}
So the initial allocation is Pareto improved upon.

Hence, once transfers are allowed, Pareto efficiency requires the ticket to go to the person with the highest willingness to pay.
\end{examplebox}

\subsection{Utilitarian Allocations}

Since transfers only change the distribution of utility, the relevant measure of the size of the ``utility pie'' generated by an allocation is the sum of $v_i(A)$ across individuals.

\begin{definition}{Utilitarian allocation}{ch11-utilitarian-allocation}
An allocation $A$ is \emph{utilitarian} if it maximizes
\[
\sum_{i=1}^n v_i(A)
\]
over all feasible allocations.
\end{definition}

Equivalently, society weakly prefers allocation $A$ to allocation $B$ under utilitarianism if and only if
\[
\sum_{i=1}^n v_i(A)\ge \sum_{i=1}^n v_i(B).
\]

\begin{examtip}
In this chapter, utilitarianism is applied to \emph{allocations}, not policies. Money transfers are dealt with separately through the vector $t$.
\end{examtip}

\subsection{Pareto Efficiency of Policies with Money Transfers}

We now define Pareto efficiency for policies, subject to balanced transfers.

\begin{definition}{Pareto improvement and Pareto efficiency for policies}{ch11-policy-pe}
A policy $(B,t')$ Pareto improves on policy $(A,t)$ if
\[
U_i(B,t_i') \ge U_i(A,t_i)
\quad \text{for all } i,
\]
and
\[
U_j(B,t_j') > U_j(A,t_j)
\quad \text{for some } j.
\]
A policy is \emph{Pareto efficient} if there is no other feasible policy with balanced transfers that Pareto improves on it.
\end{definition}

The key result of the lecture is that under quasi-linear utility and balanced transfers, Pareto-efficient policies and utilitarian allocations coincide.

\begin{theorem}{Utilitarianism and Pareto efficiency with money transfers}{ch11-main}
Suppose:
\begin{itemize}
    \item each individual has quasi-linear utility
    \[
    U_i(A,t_i)=v_i(A)+t_i;
    \]
    \item only balanced transfer schemes are allowed:
    \[
    \sum_{i=1}^n t_i=0.
    \]
\end{itemize}
Then a policy $(A,t)$ is Pareto efficient for some balanced transfer scheme $t$ if and only if $A$ is a utilitarian allocation.
\end{theorem}

\begin{proof}
We prove both directions by contraposition.

\paragraph{Direction 1.} If $A$ is utilitarian, then $(A,t)$ is Pareto efficient for some $t$.

It is enough to prove the contrapositive:
\[
\text{If $(A,t)$ is not Pareto efficient for any balanced $t$, then $A$ is not utilitarian.}
\]
Take any balanced $t$. If $(A,t)$ is not Pareto efficient, then there exists another balanced policy $(B,t')$ such that
\[
v_i(B)+t_i' \ge v_i(A)+t_i
\quad \text{for all } i,
\]
with strict inequality for at least one individual. Summing across all individuals gives
\[
\sum_{i=1}^n \bigl(v_i(B)+t_i'\bigr)
>
\sum_{i=1}^n \bigl(v_i(A)+t_i\bigr).
\]
Because both transfer schemes are balanced,
\[
\sum_{i=1}^n t_i'=\sum_{i=1}^n t_i=0,
\]
so
\[
\sum_{i=1}^n v_i(B) > \sum_{i=1}^n v_i(A).
\]
Therefore $A$ does not maximize total utility and is not utilitarian.

\paragraph{Direction 2.} If $(A,t)$ is Pareto efficient for some balanced $t$, then $A$ is utilitarian.

Again prove the contrapositive:
\[
\text{If $A$ is not utilitarian, then for any balanced $t$, $(A,t)$ is not Pareto efficient.}
\]
If $A$ is not utilitarian, then there exists some feasible allocation $B$ such that
\[
\sum_{i=1}^n v_i(B) > \sum_{i=1}^n v_i(A).
\]
Take any balanced transfer scheme $t$. Define a transfer vector
\[
\hat t_i = v_i(A)+t_i-v_i(B).
\]
Then for each individual $i$,
\[
U_i(B,\hat t_i)=v_i(B)+\hat t_i = v_i(A)+t_i = U_i(A,t_i).
\]
So each individual is indifferent between $(A,t)$ and $(B,\hat t)$.

Now sum these transfers:
\[
\sum_{i=1}^n \hat t_i
=
\sum_{i=1}^n v_i(A)+\sum_{i=1}^n t_i-\sum_{i=1}^n v_i(B)
=
\sum_{i=1}^n v_i(A)-\sum_{i=1}^n v_i(B)
<0,
\]
since $t$ is balanced and $B$ yields strictly higher total utility. Thus $\hat t$ runs a budget surplus. Let
\[
S = -\sum_{i=1}^n \hat t_i > 0.
\]
Redistribute this surplus equally by defining
\[
t_i'=\hat t_i + \frac{S}{n}.
\]
Then
\[
\sum_{i=1}^n t_i' = 0,
\]
so $t'$ is balanced. Also,
\[
U_i(B,t_i') = U_i(B,\hat t_i)+\frac{S}{n}
= U_i(A,t_i)+\frac{S}{n}
> U_i(A,t_i)
\]
for every individual $i$. Hence $(B,t')$ Pareto improves on $(A,t)$, so $(A,t)$ is not Pareto efficient.

This proves both directions.
\end{proof}

\begin{theorem}{Any balanced transfer supports a utilitarian allocation as Pareto efficient}{ch11-any-transfer}
Under the assumptions of Theorem \ref{thm:ch11-main}, if $A$ is utilitarian, then
\[
(A,t)
\]
is Pareto efficient for \emph{every} balanced transfer scheme $t$.
\end{theorem}

\begin{proof}
If some balanced policy $(B,t')$ Pareto improved on $(A,t)$, then by the summation argument in the proof of Theorem \ref{thm:ch11-main},
\[
\sum_{i=1}^n v_i(B) > \sum_{i=1}^n v_i(A),
\]
contradicting that $A$ is utilitarian.
\end{proof}

\subsection{Interpretation}

The theorem shows that, under the lecture's assumptions, the policy problem splits cleanly into two parts.

\begin{policynote}[title={Policy Interpretation: Efficiency then distribution}]
Under quasi-linear utility and balanced transfers:
\begin{enumerate}
    \item choose the allocation that maximizes
    \[
    \sum_i v_i(A),
    \]
    i.e. the largest utility pie;
    \item then use money transfers to determine how that pie is shared.
\end{enumerate}
So any non-utilitarian allocation is Pareto dominated by some policy built from a utilitarian allocation.
\end{policynote}

Graphically, the lecture interprets this as follows:
\begin{itemize}
    \item without transfers, different allocations generate a utility possibility frontier;
    \item with balanced transfers and quasi-linear utility, the frontier expands along lines of slope $-1$, because utility can be shifted across individuals through money;
    \item the outer frontier is attained by utilitarian allocations.
\end{itemize}

\begin{examtip}
The strongest intuition is:
\begin{quote}
Different allocations determine the \emph{size} of the utility pie. Transfers determine only the \emph{division} of that pie.
\end{quote}
Under quasi-linearity, any division of a smaller pie is Pareto dominated by some division of a larger pie.
\end{examtip}

\subsection{Worked Review Examples}

\begin{examplebox}[title={Worked Example: Which utility forms satisfy the utilitarianism argument?}]
Suppose policies involve allocations $(x_1^i,x_2^i)$ and transfers $t_i$. In which of the following cases does the lecture's utilitarianism argument apply?

\[
\text{(a)}\quad U_i = x_1^i x_2^i + t_i
\]
\[
\text{(b)}\quad U_i = x_1^i x_2^i
\]
\[
\text{(c)}\quad U_i = x_1^i x_2^i t_i
\]
\[
\text{(d)}\quad U_i = 0.5x_1^i + 0.5x_2^i + 0.5 t_i
\]

The argument requires utility to be quasi-linear in money, so that money enters additively with a constant marginal utility.

\begin{itemize}
    \item \textbf{(a)} Yes. This is of the form
    \[
    U_i = v_i(x)+t_i
    \]
    with
    \[
    v_i(x)=x_1^i x_2^i.
    \]
    \item \textbf{(b)} No. There is no monetary-transfer term.
    \item \textbf{(c)} No. Money enters multiplicatively, so utility is not quasi-linear in money.
    \item \textbf{(d)} Yes. This is still quasi-linear in money:
    \[
    U_i = v_i(x) + \alpha t_i
    \]
    with constant $\alpha=0.5>0$. Since the coefficient on money is constant and common, the same logic applies after a positive rescaling of utility.
\end{itemize}

Therefore the correct choices are
\[
\boxed{\text{(a) and (d).}}
\]
\end{examplebox}

\begin{examplebox}[title={Worked Example: Which allocations can be part of a Pareto-efficient policy?}]
Consider a two-person, one-good economy with balanced money transfers $(t_A,t_B)$ and utilities
\[
U_A=x_A+t_A,
\qquad
U_B=x_B+t_B.
\]
Suppose the feasible allocations are:
\[
\text{(a)}\ (x_A,x_B)=(4,4),
\qquad
\text{(b)}\ (x_A,x_B)=(5,3),
\qquad
\text{(c)}\ (x_A,x_B)=(3,5).
\]

Because
\[
v_A(x)=x_A,
\qquad
v_B(x)=x_B,
\]
the utilitarian sum for each allocation is
\[
x_A+x_B.
\]
For all three allocations,
\[
x_A+x_B=8.
\]
Hence all three allocations maximize the utilitarian objective over the listed feasible set.

By Theorem \ref{thm:ch11-main}, each utilitarian allocation can be part of a Pareto-efficient policy with some balanced transfers. In fact, by Theorem \ref{thm:ch11-any-transfer}, any balanced transfers work.

So the correct answer is
\[
\boxed{\text{(a), (b), and (c).}}
\]
\end{examplebox}

\subsection{Critical Assumptions and Limitations}

The lecture emphasizes two critical assumptions.

\begin{definition}{Critical assumptions behind the utilitarian result}{ch11-assumptions}
The utilitarian-efficiency argument requires:
\begin{enumerate}
    \item \textbf{Quasi-linear preferences in money:}
    \[
    U_i(A,t_i)=v_i(A)+t_i.
    \]
    \item \textbf{Correct implementation of balanced transfers:} society must be able to choose and enforce the needed transfer scheme after selecting the utilitarian allocation.
\end{enumerate}
\end{definition}

If these fail, then the theorem may no longer hold.

\begin{commonmistake}
Do not say that utilitarianism is always the best policy rule. The lecture's result is explicitly conditional:
\[
\text{quasi-linear utility} + \text{balanced transfers} + \text{correct implementation}.
\]
Without these, the equivalence between Pareto-efficient policies and utilitarian allocations can fail.
\end{commonmistake}

\begin{policynote}
Even when the utilitarian allocation is clear, society may still disagree over transfers. The theorem resolves the efficient \emph{allocation} question, but not the distributive question of how the gains from that allocation should be shared.
\end{policynote}

\subsection{Chapter Summary}

\begin{keyformula}[title={Key Formula: Lecture 11 Summary}]
\begin{align*}
A \to B &\Longleftrightarrow \neg B \to \neg A, \\
U_i(A,t_i) &= v_i(A)+t_i, \\
\sum_{i=1}^n t_i &= 0 \qquad \text{(balanced transfers)}, \\
\text{WTP}_i(A \to B) &= v_i(B)-v_i(A), \\
A \text{ utilitarian}
&\iff
\sum_{i=1}^n v_i(A) \ge \sum_{i=1}^n v_i(B)
\quad \text{for all feasible } B.
\end{align*}
Under quasi-linear utility and balanced transfers:
\[
(A,t) \text{ is Pareto efficient for some balanced } t
\iff
A \text{ is utilitarian.}
\]
Indeed, if $A$ is utilitarian, then
\[
(A,t)
\]
is Pareto efficient for every balanced transfer vector $t$.
\end{keyformula}

\begin{examtip}
The highest-yield takeaways from this chapter are:
\begin{itemize}
    \item know the difference between a conditional statement, its contrapositive, and its converse;
    \item know the difference between an allocation and a policy;
    \item know why balanced transfers are imposed;
    \item know how quasi-linear utility makes willingness to pay equal to utility differences;
    \item know the theorem:
    \[
    \text{Pareto-efficient policy with transfers}
    \iff
    \text{utilitarian allocation};
    \]
    \item know that the result depends critically on quasi-linear utility and the feasibility of implementing the right transfers.
\end{itemize}
\end{examtip}